\section{Assessment and exact scope}\label{sec:assessment}

\subsection{Answer to the research question}

Yes: there are useful, rigorous transfers. They fall into three mathematically different categories.

\begin{enumerate}
\item \textbf{Consequences of the exported result.} Taking the repository's Dirichlet-GRH theorem as an established input, classical theorems whose only additional assumption is Dirichlet GRH can be applied. Examples include GRH for abelian extensions of \(\Q\), logarithmic bounds for the least quadratic nonresidue, and Lindel\"of bounds for Dirichlet \(L\)-functions. Their logical implication is established mathematics; it is not a new method invented in this review.
\item \textbf{Reusable proof infrastructure.} The general comparison of covolumes, indices, and heights, and the Landau transform argument, have interfaces that can be separated from the particular automorphic construction. Sections~\ref{sec:certificate} and \ref{sec:subset} give two precise extracted statements and their proofs.
\item \textbf{New research beyond the theorem's scope.} The most concrete target identified here is the \(L\)-function of a nontrivial ideal-class character of \(\Q(\sqrt{-23})\). Its RH would imply RH for the cubic field generated by a root of \(X^3-X-1\). Section~\ref{sec:principal-target} reduces this target to a specific one-sided inequality for principal prime-power ideals. The required arithmetic inequality is not proved here.
\end{enumerate}

The distinction matters. A formal proof of a statement about one class of functions does not, by virtue of using sophisticated mathematics internally, prove the analogous assertion for every class of functions appearing in its proof. Conversely, the absence of a short human exposition does not negate a successful proof-kernel verification. The assessment below respects both points.

\subsection{The fixed source snapshot}

This review uses commit
\[
\text{\small\nolinkurl{72fead102f663d5c7a407964a13dc0cf9d977207}}
\]
of \code{lsernau/grh-lean}. Both editions pin Lean \(4.35.0\) and Mathlib commit
\[
\text{\small\nolinkurl{56a005855ea17e0b0958f73ec801ae140b0a1679}}.
\]
The readable edition is the main source for mathematical interpretation; the compact edition is the practical first choice for an independent rebuild. The repository credits Levent Alp\"oge and Claude for the proof, and describes the shortening as a collaboration involving Luke Sernau, colleagues at Anthropic, and Claude \cite{repo,announcement}.

The local audit found the following exact source counts. \emph{Code lines} means nonblank lines remaining after masking comments and strings; it is a lexical measure, not a measure of mathematical complexity.

\begin{center}
\begin{tabular}{lrr}
\toprule
Measure & Readable & Compact\\
\midrule
Lean files under \code{src} &803&794\\
Files under \code{Lib} &495&573\\
Files under \code{New} &306&219\\
Total source lines &477,124&112,827\\
Nonblank lines, including comments &432,527&95,030\\
Approximate nonblank code lines &299,902&76,149\\
Local modules imported transitively by \code{Main} &803&794\\
\bottomrule
\end{tabular}
\end{center}

Every source module belongs to the module-import closure of \code{Main}. This does not imply that every declaration in those modules occurs in the proof term of a final theorem. No import cycle was found, and no \code{Lib} module imports a \code{New} module. The inventories and the script that produced them are included in the archive.

\subsection{What the five theorems say}

The decisive interface is \src{Statements.lean}, which imports only Mathlib, together with \src{Main.lean}. It is useful to state the content in conventional mathematical notation.

\begin{longtable}{P{0.39\linewidth}P{0.55\linewidth}}
\toprule
Public declaration & Content\\
\midrule
\endhead
\code{riemannHypothesis} & Mathlib's ordinary RH for \(\zeta\).\\
\code{riemannZeta_eq_zero_re} & A zero of \(\zeta(s)\), excluding negative even integers and the separately excluded pole point \(1\), has real part \(1/2\).\\
\code{LFunction_ne_zero_of_half_lt_re} & For every nonzero modulus and every Dirichlet character \(\chi\), \(L(s,\chi)\ne0\) if \(\Rea s>1/2\), except for the principal-character pole at \(s=1\).\\
\code{LFunction_eq_zero_re} & For every such \(\chi\), a zero in \(0<\Rea s<1\) has real part \(1/2\).\\
\code{LFunction_eq_zero_re_of_isPrimitive} & The preceding statement restricted to primitive characters.\\
\bottomrule
\end{longtable}

There is no hidden primitive-character restriction in the stronger two Dirichlet statements. In particular, the result covers characters induced from a smaller conductor. Their additional Euler factors can have zeros on \(\Rea s=0\), so the open-strip formulation is the appropriate one.

The pole exclusions are also substantive. Lean represents functions by total maps, so a meromorphic function's conventional value at a pole must be excluded when expressing zero location. Excluding \(s=1\) for a principal character is correct, and does not leave a gap at an ordinary point of its analytic domain.

The exact official Mathlib definitions at the pinned commit were separately retrieved. The zeta function and Dirichlet \(L\)-function are the existing analytic continuations, agreeing with their Dirichlet series in \(\Rea s>1\); they are not newly invented functions named after those objects. This checks the outer statement boundary. It is not a reproof of all of Mathlib.

\subsection{Validation evidence and its limits}

The primary independent verification page by Andrew Sutherland is stamped 26 November 2026, 02:50 UTC, which is 25 November in Pacific time. It reports that both editions completed paranoid Comparator checking, with statement comparison, the axiom audit, Lean kernel checks, and the supplied external checkers. A second compact run also passed. A smaller-machine readable run was stopped for memory use; its lack of a verdict is separate from the successful larger-machine run \cite{checks}.

This review retrieved the relevant logs and checked source identity locally:

\begin{itemize}
\item All 1,625 entries in the verifier's source-file checksum manifest match the checkout. The only tracked files outside that list are the repository-root license and README.
\item The SHA-256 of the local Git archive matches the verifier's reported archive hash:
\begin{center}\small
\nolinkurl{9138cf2f5862e64f823ed676e8d37333c0717e1a13cd8f74f6b735ab2d1ba57c}.
\end{center}
\item All 23 selected downloaded evidence files covered by the verifier's checksum list match that list.
\end{itemize}

These matches substantiate that the source examined here is the source named in the successful verification evidence. The checksums and logs share a provenance; checking them is not the same operation as rerunning the proof kernels.

The local lexical audit found no proof placeholders or custom axioms in either \code{src} tree, and no occurrences of the execution or elaboration mechanisms included in the audit's watch list. Each top-level \code{Challenge.lean} does contain five \code{sorry} placeholders. They are comparison targets, outside the source import closure and default build targets; treating them as gaps in \code{Main} would be an error. The downloaded axiom outputs name only \code{propext}, \code{Classical.choice}, and \code{Quot.sound}.

\statusbox{\textbf{Local verification status.} Source inspection, module inventory, exact hash matching, and the arithmetic experiments in this archive were executed. Lean, Lake, and independent proof checkers were not installed in this workspace, so no local compilation or proof replay is claimed. The additional Lean interfaces have not inherited the upstream verification status.}

The repository can therefore be treated as a substantial formally checked artifact for its actual statements, supported by published independent evidence. Understanding how its construction behaves under a change of arithmetic setting remains a separate mathematical task.
